% =====================================================================
%  1. Problem Definition
% =====================================================================
\section{Problem Definition}
\label{sec:problem}

\subsection{Task}
\label{sec:task}

We forecast the hourly electricity load of Spain one day ahead. At 10:00 local time on day $D-1$,
two hours before the 12:00 gate closure of the day-ahead market, the model outputs 24 values
$\hat y_{D,0},\dots,\hat y_{D,23}$ in megawatts (MW), one for each local hour of day $D$. Only information
published before 10:00 on $D-1$ may be used. The newest load hour that is surely available at that
time is the hour stamped 08:00 on $D-1$. The most recent \emph{same-hour} load value is therefore
two days old ($D-2$), not one.

\subsection{Who needs it and what we ask}
\label{sec:questions}

The dataset already contains the forecast of the transmission system operator (TSO). If a user only
needs a number, the TSO column is the right choice, and our model does not try to replace it.
Our target user is an analyst who only has public data and wants an open, explainable model. Such a
user needs to know where to spend effort: on feature engineering, on a more complex model, or on
obtaining weather forecasts. We therefore ask three questions:
\begin{enumerate}[label=\textbf{Q\arabic*}, leftmargin=*]
  \item How much error do lag, calendar, holiday and temperature features remove, using a linear model
        (B1 Ridge versus the seasonal-naive B0)?
  \item Does a non-linear model (M1 LightGBM, Light Gradient Boosting Machine) add accuracy once these features exist (M1 versus B1)?
  \item How much is knowing the actual weather of day $D$ worth, compared with only past weather?
\end{enumerate}

\subsection{Dataset}
\label{sec:dataset}

We use the Kaggle dataset ``Hourly energy demand generation and weather''~\cite{jhana2019energy}.
It joins load and generation data from ENTSO-E (European Network of Transmission System Operators
for Electricity) with hourly weather of five Spanish cities
(\cref{tab:dataset}).

\begin{table}[H]
  \centering
  \small
  \caption{The two files of the dataset and what we use from them.}
  \label{tab:dataset}
  \begin{tabularx}{\linewidth}{@{}l l Y@{}}
    \toprule
    File & Size & Used columns \\
    \midrule
    \code{energy\_dataset.csv} & 35{,}064 hourly rows, 2015--2018, 29 columns
      & \code{total load actual} (target); \code{total load forecast} (TSO forecast, reference only).
        Generation and price columns are dropped because they are simultaneous with the target or
        published after the 10:00 cutoff. \\
    \code{weather\_features.csv} & 178{,}396 rows, 5 cities, 17 columns
      & \code{temp} of Madrid, Barcelona, Valencia, Seville and Bilbao. \\
    \bottomrule
  \end{tabularx}
\end{table}

\subsection{Two weather settings}
\label{sec:weather-settings}

The dataset has observed weather but no weather \emph{forecasts}. We therefore run every model in
two settings and never mix them inside one comparison:
\begin{itemize}
  \item \textbf{actual}: the observed temperature of day $D$ is used, as if a perfect weather
        forecast were available. This gives an optimistic bound.
  \item \textbf{lagged}: only weather observed up to 08:00 on $D-1$ is used. This is a genuine
        ex-ante forecast with lagged predictors~\cite{hyndman2021fpp3} and needs no weather feed.
        It gives a pessimistic bound.
\end{itemize}
Performance with a real weather forecast lies between the two bounds, but it cannot be measured with
this dataset.
